\chapter{Évaluation Expérimentale, Résultats et Perspectives}

Les chapitres précédents ont présenté la problématique, la méthodologie et la plateforme. Ce dernier chapitre évalue rigoureusement la solution développée : il décrit le protocole d'évaluation fondé sur un jeu de données synthétique avec vérité terrain, présente les résultats obtenus et les analyses qui les éclairent, rapporte la validation manuelle effectuée sur des cas représentatifs, et discute les limites du travail ainsi que ses perspectives.

\section{Protocole d'évaluation}

\subsection{La nécessité d'une vérité terrain}

Évaluer un système d'Entity Resolution exige de connaître, pour un ensemble de paires d'enregistrements, lesquelles correspondent réellement à une même entité. Sur les données réelles de l'institution, cette vérité terrain n'existe pas : elle nécessiterait l'annotation manuelle de milliers de paires par des experts métier, ce qui représente un effort considérable. Pour cette raison, l'évaluation repose sur un jeu de données synthétique dont la structure reproduit fidèlement celle de la table formule\_P, avec l'avantage décisif que l'entité réelle de chaque enregistrement est connue par construction.

\subsection{Construction du jeu de données synthétique}

Le jeu de données comprend cent mille déclarations réparties entre dix mille opérateurs, chaque opérateur étant représenté par un nombre variable de déclarations. Les champs générés reprennent les colonnes utilisées par le pipeline : identifiants de Registre de Commerce et leurs versions corrigées, raisons sociales déclarées et corrigées, donneur d'ordre, codes de banque et de guichet, ainsi que des champs contextuels tels que la ville et le centre.

Les variations introduites dans les données reproduisent de manière générale les phénomènes observés dans la base réelle : erreurs de frappe par substitution, insertion ou suppression de caractères ; formats hétérogènes d'un même identifiant selon l'émetteur de la déclaration ; champs laissés sans valeur sur certaines lignes ; noms d'entreprises abrégés, tronqués ou formulés différemment ; coexistence de versions déclarées et corrigées parfois divergentes. Le jeu comprend également des opérateurs proches mais distincts, destinés à vérifier que le système ne fusionne pas des entités différentes sous prétexte de leur ressemblance. La correspondance entre chaque enregistrement et son opérateur réel est conservée dans un fichier de référence qui n'est utilisé qu'au moment de l'évaluation, jamais pendant l'exécution du pipeline.

\subsection{Métriques d'évaluation}

L'évaluation utilise les métriques classiques de la classification. La précision mesure la proportion de paires correctement appariées parmi l'ensemble des paires retenues par le système : elle reflète la prudence du système. Le rappel mesure la proportion de paires réellement appariées que le système a su retrouver : il reflète son exhaustivité. Le F-score est la moyenne harmonique de la précision et du rappel, et constitue la mesure synthétique principale car elle pénalise les systèmes qui sacrifieraient entièrement l'une des deux dimensions au profit de l'autre.

\section{Résultats principaux}

Évalué sur le jeu synthétique complet avec la configuration standard du pipeline, c'est-à-dire l'entraînement EM non supervisé, le seuil de décision fixé à 0,80 et le clustering sur graphe, le système retrouve la quasi-totalité des paires réellement appariées tout en commettant très peu de fusions abusives. Autrement dit, la précision et le rappel atteignent tous deux des niveaux élevés, ce qui autorise un usage opérationnel où seules les entités de confiance intermédiaire nécessitent une révision humaine, les autres pouvant être exploitées directement. Les valeurs exactes des métriques dépendent du paramétrage retenu et de la composition du jeu d'évaluation, et ne sont donc pas présentées ici comme des garanties absolues : l'enseignement principal est la robustesse du comportement sur l'ensemble du jeu, et non un chiffre isolé.

\section{Analyses complémentaires}

\subsection{Le rôle du seuil de décision}

Le seuil appliqué à la probabilité de correspondance constitue le principal levier de réglage entre la précision et le rappel. L'analyse de la distribution des scores montre que les probabilités produites par le modèle sont fortement polarisées : les paires manifestement appariées obtiennent des scores très élevés, les paires manifestement distinctes des scores très faibles, et seule une minorité de paires ambiguës se situe dans la zone intermédiaire. Cette polarisation explique la stabilité des résultats autour du seuil retenu de 0,80 : déplacer légèrement le seuil ne modifie qu'à la marge l'ensemble des paires retenues, car peu de paires se trouvent au voisinage immédiat du seuil. Relever fortement le seuil améliorerait marginalement la précision au prix d'une perte de rappel sur les cas ambigus, tandis que l'abaisser produirait l'effet inverse.

\subsection{L'importance de la normalisation}

Les analyses d'ablation, qui consistent à désactiver un composant du pipeline pour mesurer sa contribution, montrent que la normalisation des champs joue un rôle déterminant. Lorsque les identifiants sont comparés sous leurs formes brutes, les variations de format masquent les ressemblances et le modèle ne peut pas exploiter pleinement les indices disponibles. Après normalisation, les formes superficielles disparaissent et les fonctions de similarité se concentrent sur les différences significatives. Cet effet est particulièrement marqué sur les identifiants numériques sujets aux variations de présentation et sur les noms d'entreprises sujets aux variations de formulation des formes juridiques. La leçon générale est que la qualité de la préparation des données conditionne davantage les performances finales que la sophistication du modèle de scoring.

\subsection{Le blocking comme étape critique}

Le blocking détermine l'ensemble des paires que le modèle aura l'occasion d'examiner, et constitue donc un plafond pour le rappel du système. Une stratégie de blocking trop restrictive élimine définitivement des paires réellement appariées, qu'aucune étape ultérieure ne pourra récupérer. À l'inverse, une stratégie trop large noie le modèle sous des candidats non pertinents et alourdit considérablement le calcul. La stratégie multi-passe retenue équilibre ces deux exigences en combinant des règles strictes qui capturent efficacement les cas évidents et des règles larges qui servent de filet de sécurité pour les cas dégradés. L'expérience confirme que le soin apporté aux règles de blocage est au moins aussi important que le réglage du modèle lui-même.

\subsection{Apport des embeddings neuronaux}

La comparaison entre les approches de similarité textuelle montre un résultat instructif : sur les noms d'entreprises, une représentation classique fondée sur les termes discriminants obtient une précision légèrement supérieure à celle des plongements neuronaux multilingues. Ce résultat s'explique par la nature des données : les noms d'entreprises sont des séquences courtes où quelques termes distinctifs suffisent à discriminer les entités, et où la compréhension fine du sens apporte peu par rapport à la reconnaissance des termes. Les plongements conservent néanmoins un intérêt pour les cas ambigus où les formulations diffèrent profondément tout en désignant la même entreprise, et leur combinaison à parts égales avec le score probabiliste n'altère pas les performances globales. Ils constituent donc un enrichissement utile mais non décisif sur ce type de données.

\subsection{Comparaison des stratégies de clustering}

Les deux stratégies de clustering ont été comparées sur les mêmes paires retenues. Les composantes connexes produisent des groupes complets mais exposés au phénomène d'enchaînement, où un lien faible suffit à fusionner artificiellement deux entités distinctes. L'algorithme de Leiden produit des groupes plus resserrés en ne conservant que les communautés à structure interne forte, au prix d'un temps de calcul légèrement supérieur qui reste néanmoins négligeable devant celui de l'entraînement et de la prédiction. Pour un usage opérationnel où les fusions abusives sont plus coûteuses que les fragmentations résiduelles, la sélectivité de Leiden constitue un avantage.

\subsection{Synthèse qualitative des analyses}

Le tableau suivant résume l'apport observé de chaque composant, sans valeur chiffrée, tel qu'il ressort des analyses d'ablation et des comparaisons menées pendant le projet.

\begin{table}[H]
\centering
\caption{Synthèse qualitative de l'apport de chaque composant du pipeline.}
\renewcommand{\arraystretch}{1.3}
\setlength{\tabcolsep}{3pt}
{\small
\begin{tabular}{|p{3.6cm}|p{10.0cm}|}
\hline
\textbf{Composant} & \textbf{Effet observé} \\
\hline
Normalisation des champs & Amélioration déterminante : les variations de format masquaient les ressemblances avant normalisation \\
\hline
Blocking multi-passe & Plafonne le rappel du système : les règles larges servent de filet de sécurité pour les cas dégradés \\
\hline
Seuil de décision & Comportement stable autour du seuil retenu, car les scores sont fortement polarisés \\
\hline
Embeddings neuronaux & Enrichissement utile pour les cas ambigus, sans effet décisif sur les noms courts et standardisés \\
\hline
Clustering de Leiden & Groupes plus sélectifs que les composantes connexes, préférables quand les fusions abusives sont coûteuses \\
\hline
\end{tabular}
}
\end{table}

\section{Validation manuelle}

\subsection{Protocole de validation}

En complément de l'évaluation quantitative, un échantillon de résultats a été examiné manuellement pour apprécier la qualité des entités produites dans des situations représentatives. Trois catégories de cas ont été considérées : les doublons manifestes, où des enregistrements aux noms quasi identiques doivent être regroupés ; les quasi-doublons trompeurs, où des enregistrements aux noms similaires mais aux identifiants divergents doivent rester séparés ; et les cas limites, où la décision nécessite un jugement attentif sur l'ensemble des indices disponibles.

\subsection{Observations}

Sur les cas manifestes, le système se comporte comme attendu dans la quasi-totalité des situations. Sur les quasi-doublons trompeurs, la combinaison des indices permet généralement de maintenir la séparation des entités distinctes, notamment grâce au poids accordé aux identifiants structurés. Les cas limites restants relèvent principalement de deux situations : les entreprises homonymes, c'est-à-dire des entreprises distinctes qui partagent exactement le même nom et parfois la même activité, que seul un examen des identifiants officiels permet de départager ; et les entreprises existant sous plusieurs formes juridiques successives, où la décision de fusionner ou non dépend d'une connaissance de l'histoire de l'entreprise que les données ne contiennent pas toujours. Ces situations justifient le maintien d'une révision humaine pour les entités de confiance intermédiaire. Le tableau suivant résume les cas examinés et le comportement attendu du système pour chacun.

\begin{table}[H]
\centering
\caption{Cas de validation manuelle et comportement attendu du système.}
\renewcommand{\arraystretch}{1.3}
\setlength{\tabcolsep}{3pt}
{\small
\begin{tabular}{|p{4.5cm}|p{9.1cm}|}
\hline
\textbf{Cas} & \textbf{Comportement attendu} \\
\hline
Doublons manifestes & Regroupement des enregistrements aux noms quasi identiques en une seule entité \\
\hline
Quasi-doublons trompeurs & Maintien de la séparation des entités distinctes malgré la ressemblance des noms \\
\hline
Entreprises homonymes & Révision humaine, car le nom seul ne permet pas de trancher \\
\hline
Formes juridiques successives & Révision humaine, car la décision dépend de l'histoire de l'entreprise \\
\hline
\end{tabular}
}
\end{table}

\section{Limites}

Trois limites principales doivent être signalées :

\begin{itemize}
    \item La nature synthétique des données d'évaluation : aussi fidèle soit la reproduction des phénomènes observés, un jeu synthétique ne capture pas nécessairement toutes les spécificités des données historiques accumulées sur plusieurs années, avec leurs évolutions de format et leurs cas particuliers. Une validation complémentaire sur un échantillon de données réelles annotées manuellement reste donc nécessaire avant tout usage en production.
    \item Le passage à l'échelle : le pipeline a été évalué sur des volumes de l'ordre de la centaine de milliers de déclarations, tandis que la base réelle en contient plusieurs millions. Les étapes les plus coûteuses, à savoir le blocking, l'entraînement EM et la prédiction, devront être testées et éventuellement optimisées sur ces volumes, par exemple par un traitement partitionné ou par une parallélisation accrue.
    \item L'hypothèse d'indépendance conditionnelle des champs sur laquelle repose le modèle probabiliste : en pratique, certains champs sont corrélés entre eux, comme le code de guichet qui dépend de la banque, ou la forme juridique qui fait partie du nom. Cette hypothèse simplificatrice n'empêche pas d'excellentes performances empiriques, mais elle peut biaiser les probabilités estimées dans les cas ambigus.
\end{itemize}

\section{Perspectives}

Les perspectives s'ordonnent en trois horizons. À court terme, la priorité est la validation du pipeline sur un échantillon de données réelles avec revue manuelle des entités par les experts métier de l'institution. Cette validation permettra de confirmer les performances mesurées sur données synthétiques et d'identifier les adaptations nécessaires aux spécificités des données historiques.

À moyen terme, plusieurs enrichissements sont envisageables :

\begin{itemize}
    \item L'optimisation des performances pour le traitement de volumes de plusieurs millions de déclarations.
    \item L'intégration du pipeline dans le processus périodique de chargement des déclarations, afin que les nouvelles déclarations soient rattachées aux entités existantes au fil de l'eau.
    \item Le développement d'un tableau de bord de suivi de la qualité des identifiants, qui alerterait sur les dégradations et orienterait les efforts de fiabilisation à la source.
    \item La mise en place d'un apprentissage actif, où les corrections apportées par les analystes lors de la révision alimenteraient l'amélioration continue du système.
\end{itemize}

À long terme, le périmètre pourrait être étendu à d'autres tables de déclarations, et les entités résolues pourraient être rapprochées des référentiels externes d'entreprises pour enrichir encore la connaissance des opérateurs.

\section*{Conclusion}

Ce chapitre a démontré empiriquement l'efficacité de l'approche retenue : sur un jeu de données représentatif avec vérité terrain, le système retrouve la quasi-totalité des paires réellement appariées tout en commettant très peu de fusions abusives, et ce sans aucune donnée annotée grâce à l'entraînement non supervisé. Les analyses ont éclairé le rôle de chaque composant et confirmé que la préparation des données et le blocking sont au moins aussi déterminants que le modèle de scoring. La validation manuelle a circonscrit les situations qui justifient le maintien d'une révision humaine. Les limites identifiées tracent une feuille de route claire vers un usage en production, dont la première étape est la validation sur données réelles.
